% Define colours for code listings
\definecolor{keyword}{RGB}{0,0,128}
\definecolor{comment}{RGB}{0,96,0}
\definecolor{string}{RGB}{128,0,0}
\lstset{
  language=Python,
  basicstyle=\ttfamily\small,
  keywordstyle=\color{keyword}\bfseries,
  commentstyle=\color{comment}\itshape,
  stringstyle=\color{string},
  showstringspaces=false,
  frame=single,
  framerule=0.5pt,
  numbers=left,
  numberstyle=\tiny,
  numbersep=5pt
}

\chapter{Lab 6: Building the Customer‑Facing Frontend}
\label{chap:Lab 6: Building the Customer‑Facing Frontend}

\section{Introduction and Objectives}

With your customer‑support agent now deployed to a production runtime
and continuously evaluated for quality, the final piece of the
workshop is to provide a \emph{customer interface}.  Lab 6 guides you
through building a Streamlit web application that allows end users to
chat with your agent in a secure, user‑friendly environment.  Key
features of the frontend include:

\begin{itemize}
  \item \textbf{Secure Authentication} – Users log in via Amazon Cognito,
        ensuring that only authorised customers can access the agent.
  \item \textbf{Real‑time Chat Interface} – A conversational UI built
        with Streamlit that resembles a modern messaging application.
  \item \textbf{Streaming Responses} – Live response streaming for
        improved user experience; messages appear as they are
        generated rather than waiting for the full answer.
  \item \textbf{Session Management} – The frontend passes
        session IDs to AgentCore Runtime and maintains conversation
        history between page refreshes, leveraging the memory
        introduced in Lab 2.
  \item \textbf{Performance Metrics} – Each response displays a timing
        indicator so customers see how quickly the agent handled their
        query.
\end{itemize}

By the end of this lab you will have:

\begin{itemize}
  \item Deployed a Streamlit application connected to your AgentCore
        Runtime endpoint.
  \item Integrated Cognito authentication into the frontend.
  \item Implemented real‑time streaming of agent responses.
  \item Completed a fully self‑contained customer support solution
        spanning prototype through production.
\end{itemize}

\section{Prerequisites}

Before beginning the lab make sure you meet the following conditions:

\begin{itemize}
  \item \textbf{Completed Labs 1–4.}  You need a deployed AgentCore
        Runtime endpoint and a configured Cognito user pool from
        earlier labs.
  \item \textbf{Python 3.10 or higher} installed on your machine.
  \item \textbf{Streamlit dependencies.}  The lab installs these via
        the provided requirements file.
  \item \textbf{AgentCore Runtime ARN and Cognito credentials.}
        Ensure you know the ARN of the runtime and have test user
        credentials for Cognito authentication.
\end{itemize}

\section{Installing Frontend Dependencies}

The first step is to install the packages needed by the Streamlit
frontend.  These include Streamlit itself, the Cognito
authentication library and helper modules for interacting with AgentCore.
Run the following command from the root of the workshop repository:

\begin{lstlisting}[language=bash]
pip install -r lab_helpers/lab5_frontend/requirements.txt -q
echo "✅ Frontend dependencies installed successfully!"
\end{lstlisting}

This command reads the \texttt{requirements.txt} file in the
\texttt{lab\_helpers/lab5\_frontend} directory and installs all necessary
packages, suppressing verbose output for a cleaner notebook.

\section{Frontend Architecture}

The Streamlit application is split into a few key files.  Each file
encapsulates a distinct responsibility in the frontend and allows
you to customise or extend the interface without touching the core
agent logic.

\subsection{Core Components}

\begin{description}
  \item[\texttt{main.py}] This is the entry point for the Streamlit app.  It
    handles user authentication using a \texttt{CognitoAuthenticator},
    displays the chat interface, maintains session state and
    orchestrates streaming of responses.  After the user logs in,
    \texttt{main.py} creates a \texttt{ChatManager} instance to
    communicate with AgentCore Runtime and renders each message in the
    conversation.

  \item[\texttt{chat.py}] Implements the \texttt{ChatManager} class and
    functions for invoking the AgentCore Runtime.  It offers both
    streaming and non‑streaming invocation methods.  The streaming
    variant uses the HTTP \texttt{text/event-stream} mechanism to
    yield partial results as the model generates them.

  \item[\texttt{chat\_utils.py}] Provides helper functions for reading and
    writing configuration from Systems Manager Parameter Store,
    converting plain text into clickable links, reading YAML/JSON
    configuration files and creating safe HTML for Streamlit.

  \item[\texttt{sagemaker\_helper.py}] Supplies a helper function to retrieve
    the URL of a SageMaker Studio Lab environment.  If the code is not
    running in SageMaker, it falls back to \texttt{localhost:8501}.
\end{description}

\subsection{Understanding \texttt{main.py}}

The heart of the frontend is the \texttt{main.py} script.  It begins
by loading a JSON secret from a local helper function and instantiating
a \texttt{CognitoAuthenticator}.  The authenticator prompts the user
for their username and password; if authentication fails the
application halts.

\begin{lstlisting}
import json
import os
import sys
import time
import uuid

import streamlit as st
from chat import ChatManager, invoke_endpoint_streaming
from chat_utils import make_urls_clickable
from streamlit_cognito_auth import CognitoAuthenticator

from utils import get_customer_support_secret

secret = json.loads(get_customer_support_secret())

authenticator = CognitoAuthenticator(
    pool_id=secret["pool_id"],
    app_client_id=secret["client_id"],
    app_client_secret=secret["client_secret"],
    use_cookies=False,
)

is_logged_in = authenticator.login()
if not is_logged_in:
    st.stop()

# Once authenticated, display the app title and set up sidebar
st.title("Customer Support Agent")
with st.sidebar:
    st.text(f"Welcome,\n{authenticator.get_username()}")
    st.button("Logout", on_click=authenticator.logout)

\end{lstlisting}

After successful login, the script initialises a \texttt{ChatManager}
and session state variables.  Streamlit’s session state persists
across page reloads, enabling conversation history and session
identification.  Each time the user submits a prompt, \texttt{main.py}
appends it to the session message list and then streams the
assistant’s response.  The following code excerpt shows how a
user message is processed:

\begin{lstlisting}
# Accept user input
if prompt := st.chat_input("What is up?"):
    # Display user message
    with st.chat_message("user"):
        st.markdown(prompt)
    # Add user message to history
    st.session_state.messages.append({"role": "user", "content": prompt})
    # Build a conversation context window
    context = "".join(
        f"{msg['role'].capitalize()}: {msg['content']}\n"
        for msg in st.session_state.messages[-20:]
    )
    # Retrieve the bearer token for the logged‑in user
    bearer_token = st.session_state.get("auth_access_token")
    # Stream the response
    with st.chat_message("assistant"):
        message_placeholder = st.empty()
        start_time = time.time()
        accumulated = ""
        for chunk in invoke_endpoint_streaming(
            agent_arn=st.session_state["agent_arn"],
            payload=json.dumps({"prompt": context}),
            session_id=st.session_state["session_id"],
            bearer_token=bearer_token,
        ):
            accumulated += chunk
            clickable = make_urls_clickable(accumulated)
            message_placeholder.markdown(
                f"<div class='assistant-bubble streaming'> {clickable}</div>",
                unsafe_allow_html=True,
            )
        elapsed = time.time() - start_time
        # Final response rendering with timing
        message_placeholder.markdown(
            f"<div class='assistant-bubble'> {clickable}<br><span style='font-size:0.9em;color:#888;'>⏱️ Response time: {elapsed:.2f} seconds</span></div>",
            unsafe_allow_html=True,
        )
        # Persist assistant response in session history
        st.session_state.messages.append({"role": "assistant", "content": accumulated, "elapsed": elapsed})

\end{lstlisting}

This pattern demonstrates how the application accumulates chunks from
the streaming invocation, updates the placeholder to show partial
results and records the full answer along with response timing.

\subsection{Streaming Invocation with \texttt{chat.py}}

The \texttt{invoke\_endpoint\_streaming} function defined in
\texttt{chat.py} sends an HTTP POST request to AgentCore Runtime with
appropriate headers and yields chunks of the streaming response.  A
bare‑bones version is shown below:

\begin{lstlisting}
import json
import urllib.parse
import requests

def invoke_endpoint_streaming(agent_arn: str, payload: str, session_id: str, bearer_token: str, endpoint_name: str = "DEFAULT"):
    # Escape ARN and build URL
    escaped_arn = urllib.parse.quote(agent_arn, safe="")
    url = f"https://bedrock-agentcore.{st.session_state['region']}.amazonaws.com/runtimes/{escaped_arn}/invocations"
    # Build headers including the JWT token and session ID
    headers = {
        "Authorization": f"Bearer {bearer_token}",
        "Content-Type": "application/json",
        "X-Amzn-Bedrock-AgentCore-Runtime-Session-Id": session_id,
    }
    # Parse payload and make streaming request
    body = json.loads(payload) if isinstance(payload, str) else payload
    response = requests.post(url, headers=headers, params={"qualifier": endpoint_name}, json=body, stream=True, timeout=100)
    response.raise_for_status()
    # Yield each chunk prefixed with 'data: '
    for line in response.iter_lines(chunk_size=1, decode_unicode=True):
        if line and line.startswith("data: "):
            chunk = line[6:]  # Remove prefix
            if chunk.strip():
                yield chunk

\end{lstlisting}

This function abstracts away the details of streaming and allows
\texttt{main.py} to focus on user interface concerns.  A similar
\texttt{invoke\_endpoint\_nostreaming} method is available for
synchronous requests.

\section{Launching the Application}

Once you understand the architecture, you can start the frontend and
connect it to your runtime.  The following commands (also found in
the notebook) first retrieve an accessible URL for the application and
then launch the Streamlit server:

\begin{lstlisting}
from lab_helpers.lab5_frontend.sagemaker_helper import get_streamlit_url

streamlit_url = get_streamlit_url()
print(f"\n Customer Support Streamlit Application URL:\n{streamlit_url}\n")

# Start the application.  This runs until you interrupt it (Ctrl+C).
!cd lab_helpers/lab5_frontend/ && streamlit run main.py
\end{lstlisting}

If you are running inside SageMaker Studio Lab, the helper extracts the
domain ID and space name from the environment and constructs a
publicly accessible URL.  Otherwise it falls back to
\texttt{http://localhost:8501}.  When the command completes, open
the provided URL in your browser and log in using your Cognito
credentials.

\section{Testing the Customer Support Application}

With the frontend running, you can simulate real customer sessions.
After logging in, you should see a welcome message followed by a chat
prompt.  Try the following scenarios:

\begin{enumerate}
  \item \textbf{Product Information Queries} – Ask about the
    specifications of your company’s products (e.g. “What are the
    specifications for your laptops?”).  The agent should call the
    product information tool and respond with the appropriate details.

  \item \textbf{Return Policy Questions} – Ask about the return policy
    for various product categories (e.g. “What’s the return policy for
    electronics?”).  The agent should use the return‑policy tool.

  \item \textbf{Troubleshooting Support} – Describe a technical
    problem (e.g. “My iPhone is overheating, what should I do?”).  This
    triggers the technical‑support tool, which queries the knowledge
    base and returns step‑by‑step instructions.

  \item \textbf{Memory and Personalisation} – Conduct a multi‑turn
    conversation, refresh the page and then continue speaking.  The
    agent should remember your prior preferences and queries thanks to
    the memory service from Lab 2.
\end{enumerate}

As you test, pay attention to:

\begin{itemize}
  \item \textbf{Real‑time Streaming} – Responses should arrive
        incrementally with a typing animation.
  \item \textbf{Response Timing} – A stopwatch icon displays how long
        each response took to generate.
  \item \textbf{Tool Integration} – Check that the agent selects the
        correct tools for different query types.
  \item \textbf{Memory Persistence} – Verify that conversation
        context persists across browser refreshes and logins.
  \item \textbf{Error Handling} – Observe how the interface responds
        if the agent encounters an error; appropriate messages should
        be displayed rather than raw tracebacks.
\end{itemize}

\section{Summary and Next Steps}

In this lab you built and tested a customer‑facing Streamlit web
application that integrates securely with your AgentCore Runtime.  You
learned how to:

\begin{itemize}
  \item Install Streamlit and helper dependencies.
  \item Authenticate users via Cognito and manage JWT tokens in the
        frontend.
  \item Stream responses from AgentCore Runtime and present them in
        an interactive chat interface with response timing.
  \item Maintain session state and conversation history using
        Streamlit’s session object.
  \item Test various customer support scenarios and verify that the
        agent uses the right tools while remembering prior context.
\end{itemize}

With Lab 6 complete, you now have a full end‑to‑end solution: a
customer‑support agent defined in Lab 1, enriched with memory in
Lab 2, scaled via Gateway and identity in Lab 3, deployed to a
production runtime in Lab 4, continuously evaluated in Lab 5 and
presented to customers through a polished web interface in Lab 6.

Future enhancements might include custom styling, integration with
existing CRM systems, multi‑language support or a mobile‑friendly
layout.  The modular architecture of the frontend makes these
improvements straightforward to implement.

\end{document}